% GENERATED FILE. Edit canonical lessons, then run npm run generate:books.
% canonical-source-hash: fnv1a64:f2f9d393f0930060
% canonical-lessons: ZH-C100-yinhang, ZH-C100-chaoshi, ZH-C100-gongyuan, ZH-C100-tushuguan, ZH-C100-youju

\chapter{Around Town}
\label{ch:zh-around-town}
\hlchaptermodality{100}

\begin{chapteropening}
\textbf{By the end of this chapter:} \emph{I can name the bank, the supermarket, the park, the library and the post office.}

Complete the last lesson of chapter 100: I can name the bank, the supermarket, the park, the library and the post office.
\end{chapteropening}

\section[yínháng]{yínháng — bank}

\label{lesson:ZH-C100-yinhang}



\begin{quote}
Before the new one: say the Mandarin for to have a shower, then the Mandarin for to comb one's hair.
\end{quote}

\subsection*{You'll want to know: yínháng}
\emph{yínháng} — \textquotedblleft{}bank\textquotedblright{}.

\begin{quote}
\emph{yínháng zài nǎr} — Where's the bank?
\end{quote}

Its characters come later; for now it is a word for the ear.

\begin{cousinweb}
\emph{yínháng} is built from yín 'silver' and háng 'a row', both inherited from Middle Chinese, \emph{ngin} and \emph{hang} (Guangyun, in Baxter's transcription). Háng is usually explained as a row of shops and then a trade; the step is not given there, so how the pair came to mean a bank is not settled.
\end{cousinweb}

\begin{grammarlens}[title={What you've built}]
One more word for asking where things are.
\end{grammarlens}

\subsection*{Your turn}
\begin{itemize}
\raggedright
  \item \emph{Say it:} \emph{yínháng}
  \item \emph{Say it:} \emph{yínháng}, once more
  \item \emph{Say it:} \emph{yínháng zài nǎr} — Where's the bank?
  \item \emph{Recall:} say the Mandarin for to have a shower, then the Mandarin for to comb one's hair, then say \emph{yínháng} again
\end{itemize}

\begin{tcolorbox}[breakable,colback=teal!4,colframe=teal!35!black,title={Before you move on}]
What does yínháng mean? (\textquotedblleft{}Bank\textquotedblright{}.) Say it once more.
\end{tcolorbox}
\section[chāoshì]{chāoshì — supermarket}

\label{lesson:ZH-C100-chaoshi}



\begin{quote}
Before the new one: say the Mandarin for to comb one's hair, then the Mandarin for bank.
\end{quote}

\subsection*{You'll want to know: chāoshì}
\emph{chāoshì} — \textquotedblleft{}supermarket\textquotedblright{}.

\begin{quote}
\emph{wǒ qù chāoshì mǎi shuǐguǒ} — I'm off to the supermarket for some fruit.
\end{quote}

Its characters come later; for now it is a word for the ear.

\begin{cousinweb}
\emph{chāoshì} is built from chāo and shì 'market', both inherited from Middle Chinese, \emph{trhjew} and \emph{dzyiX} (Guangyun, in Baxter's transcription). The dictionary glosses the first 'to leap'; 'super' is its sense in this modern word.
\end{cousinweb}

\begin{grammarlens}[title={What you've built}]
One more word for asking where things are.
\end{grammarlens}

\subsection*{Your turn}
\begin{itemize}
\raggedright
  \item \emph{Say it:} \emph{chāoshì}
  \item \emph{Say it:} \emph{chāoshì}, once more
  \item \emph{Say it:} \emph{wǒ qù chāoshì mǎi shuǐguǒ} — I'm off to the supermarket for some fruit.
  \item \emph{Recall:} say the Mandarin for to comb one's hair, then the Mandarin for bank, then say \emph{chāoshì} again
\end{itemize}

\begin{tcolorbox}[breakable,colback=teal!4,colframe=teal!35!black,title={Before you move on}]
What does chāoshì mean? (\textquotedblleft{}Supermarket\textquotedblright{}.) Say it once more.
\end{tcolorbox}
\section[gōngyuán]{gōngyuán — park}

\label{lesson:ZH-C100-gongyuan}



\begin{quote}
Before the new one: say the Mandarin for bank, then the Mandarin for supermarket.
\end{quote}

\subsection*{You'll want to know: gōngyuán}
\emph{gōngyuán} — \textquotedblleft{}park\textquotedblright{}.

\begin{quote}
\emph{wǒmen qù gōngyuán wán} — We're going to the park to play.
\end{quote}

Its characters come later; for now it is a word for the ear.

\begin{cousinweb}
\emph{gōngyuán} is built from gōng 'public' and yuán 'garden', both inherited from Middle Chinese, \emph{kuwng} and \emph{hjwon} (Guangyun, in Baxter's transcription): a garden for everyone.
\end{cousinweb}

\begin{grammarlens}[title={What you've built}]
One more word for asking where things are.
\end{grammarlens}

\subsection*{Your turn}
\begin{itemize}
\raggedright
  \item \emph{Say it:} \emph{gōngyuán}
  \item \emph{Say it:} \emph{gōngyuán}, once more
  \item \emph{Say it:} \emph{wǒmen qù gōngyuán wán} — We're going to the park to play.
  \item \emph{Recall:} say the Mandarin for bank, then the Mandarin for supermarket, then say \emph{gōngyuán} again
\end{itemize}

\begin{tcolorbox}[breakable,colback=teal!4,colframe=teal!35!black,title={Before you move on}]
What does gōngyuán mean? (\textquotedblleft{}Park\textquotedblright{}.) Say it once more.
\end{tcolorbox}
\section[túshūguǎn]{túshūguǎn — library}

\label{lesson:ZH-C100-tushuguan}



\begin{quote}
Before the new one: say the Mandarin for supermarket, then the Mandarin for park.
\end{quote}

\subsection*{You'll want to know: túshūguǎn}
\emph{túshūguǎn} — \textquotedblleft{}library\textquotedblright{}.

\begin{quote}
\emph{wǒ qù túshūguǎn kàn shū} — I'm going to the library to read.
\end{quote}

Its characters come later; for now it is a word for the ear.

\begin{cousinweb}
\emph{túshūguǎn} is built from tú 'picture', \zh{书} 'book' and guǎn 'hall', all inherited from Middle Chinese, \emph{du}, \emph{syo} and \emph{kwanH} (Guangyun, in Baxter's transcription); the dictionary glosses the first 'to plan' and the last as a lodging house for guests.
\end{cousinweb}

\begin{grammarlens}[title={What you've built}]
One more word for asking where things are.
\end{grammarlens}

\subsection*{Your turn}
\begin{itemize}
\raggedright
  \item \emph{Say it:} \emph{túshūguǎn}
  \item \emph{Say it:} \emph{túshūguǎn}, once more
  \item \emph{Say it:} \emph{wǒ qù túshūguǎn kàn shū} — I'm going to the library to read.
  \item \emph{Recall:} say the Mandarin for supermarket, then the Mandarin for park, then say \emph{túshūguǎn} again
\end{itemize}

\begin{tcolorbox}[breakable,colback=teal!4,colframe=teal!35!black,title={Before you move on}]
What does túshūguǎn mean? (\textquotedblleft{}Library\textquotedblright{}.) Say it once more.
\end{tcolorbox}
\section[yóujú]{yóujú — post office}

\label{lesson:ZH-C100-youju}



\begin{quote}
Before the new one: say the Mandarin for park, then the Mandarin for library.
\end{quote}

\subsection*{You'll want to know: yóujú}
\emph{yóujú} — \textquotedblleft{}post office\textquotedblright{}.

\begin{quote}
\emph{yóujú zài yínháng pángbiān} — The post office is next to the bank.
\end{quote}

Its characters come later; for now it is a word for the ear.

\begin{cousinweb}
\emph{yóujú} is built from yóu 'post' and jú 'office', both inherited from Middle Chinese, \emph{hjuw} and \emph{gjowk} (Guangyun, in Baxter's transcription); the dictionary glosses the first as a lodge on the border and the second as a government office.
\end{cousinweb}

\begin{grammarlens}[title={What you've built}]
One more word for asking where things are.
\end{grammarlens}

\subsection*{Your turn}
\begin{itemize}
\raggedright
  \item \emph{Say it:} \emph{yóujú}
  \item \emph{Say it:} \emph{yóujú}, once more
  \item \emph{Say it:} \emph{yóujú zài yínháng pángbiān} — The post office is next to the bank.
  \item \emph{Recall:} say the Mandarin for park, then the Mandarin for library, then say \emph{yóujú} again
\end{itemize}

\begin{tcolorbox}[breakable,colback=teal!4,colframe=teal!35!black,title={Before you move on}]
What does yóujú mean? (\textquotedblleft{}Post office\textquotedblright{}.) Say it once more.
\end{tcolorbox}
